\subsection{三角置換}

$x=\tan\theta$,\,$-\frac\pi2<\theta<\frac\pi2$と置く.この範囲では$\sec\theta>0$なので
\[
\sqrt{1+x^2}=\sec\theta,\qquad
\dd x=\sec^2\theta\dd\theta,\qquad
\int\sqrt{1+x^2}\dd x=\int\sec^3\theta\dd\theta.
\]
根号はなくなったが,三乗の積分が残った.\,$J=\int\sec^3\theta\dd\theta$と書き,\,$\sec\theta$を微分する側,\,$\sec^2\theta$を積分する側にして部分積分する.
\[
\begin{aligned}
J&=\sec\theta\tan\theta
-\int\sec\theta\tan^2\theta\dd\theta\\
&=\sec\theta\tan\theta
-\int\sec\theta(\sec^2\theta-1)\dd\theta\\
&=\sec\theta\tan\theta-J+\int\sec\theta\dd\theta.
\end{aligned}
\]
任意定数を最後にまとめるものとして,この等式を$J$について解く.前に求めた正割の積分から
\[
J=\frac12\sec\theta\tan\theta
+\frac12\log(\sec\theta+\tan\theta)+C.
\]
この範囲では$\sec\theta>|\tan\theta|$だから,対数の中は正である.\,$x$へ戻して
\begin{equation}\label{式:一と二乗の根号積分}
\int\sqrt{1+x^2}\dd x
=\frac12\left\{x\sqrt{1+x^2}
+\log\bigl(x+\sqrt{1+x^2}\bigr)\right\}+C
\end{equation}
を得る.三角置換だけで計算が終わるのではなく,その後に部分積分を使った.


